\documentclass[aspectratio=169,11pt]{beamer}

\usetheme{AnnArbor}
% \usecolortheme{seahorse}
% \usefonttheme{professionalfonts}

% \usepackage{fontspec}
% \setsansfont{DejaVu Sans}
% \setmainfont{DejaVu Serif}
% \setmonofont{DejaVu Sans Mono}
\usepackage{amsmath,amssymb,booktabs,array,multicol}
\usepackage{listings}
\usepackage{xcolor}
\usepackage{graphicx}
\usepackage{hyperref}

\definecolor{CodeBg}{RGB}{248,248,248}
\definecolor{TitleBlue}{RGB}{38,79,129}
\definecolor{AccentGreen}{RGB}{35,125,93}
\definecolor{SoftGray}{RGB}{80,80,80}
\definecolor{QuizOrange}{RGB}{180,95,30}

\setbeamercolor{title}{fg=TitleBlue}
\setbeamercolor{frametitle}{fg=TitleBlue}
\setbeamercolor{structure}{fg=AccentGreen}
\setbeamercolor{block title}{fg=white,bg=TitleBlue}
\setbeamercolor{block body}{bg=blue!3}
\setbeamercolor{alertblock title}{fg=white,bg=QuizOrange}
\setbeamercolor{alertblock body}{bg=orange!6}
\setbeamertemplate{navigation symbols}{}
\setbeamertemplate{footline}[frame number]
\setbeamertemplate{itemize item}{\raise0.15ex\hbox{\tiny$\blacktriangleright$}}
\setbeamertemplate{itemize subitem}{\raise0.15ex\hbox{\tiny$\bullet$}}

\lstdefinestyle{py}{
  language=Python,
  basicstyle=\ttfamily\scriptsize,
  backgroundcolor=\color{CodeBg},
  frame=single,
  framerule=0.2pt,
  rulecolor=\color{gray!35},
  breaklines=true,
  showstringspaces=false,
  keywordstyle=\color{TitleBlue}\bfseries,
  commentstyle=\color{SoftGray}\itshape,
  stringstyle=\color{AccentGreen}
}

\newcommand{\key}[1]{\textbf{\textcolor{AccentGreen}{#1}}}
\newcommand{\smallnote}[1]{\vspace{0.25em}{\footnotesize\color{SoftGray}#1}}
\newcommand{\quiztitle}[1]{\textbf{\textcolor{QuizOrange}{#1}}}
\newcommand{\imageplaceholder}[2]{%
\begin{center}
\IfFileExists{#1}{%
  \includegraphics[width=0.88\linewidth,height=0.62\textheight,keepaspectratio]{#1}\\[0.4em]
  {\small\textcolor{SoftGray}{\textbf{#2}}}
}{%
  \fcolorbox{SoftGray}{gray!5}{%
  \begin{minipage}[c][0.60\textheight][c]{0.84\linewidth}
  \centering
  {\Large\textcolor{TitleBlue}{Image Placeholder}}\\[1.0em]
  {\normalsize\textbf{#2}}\\[1.2em]
  {\footnotesize Source in the Markdown file:}\\[0.35em]
  {\small\texttt{\detokenize{#1}}}
  \end{minipage}}
}%
\end{center}
}

\title[Data Analysis Foundations]{Data Analysis with Python}
\subtitle{What Data Analysis Is and the Six-Step Data Analysis Process}
\author{Duc-Minh Vu}
\institute{SLSCM Lab - FDA - NEU}
\date{\today}

\begin{document}

% 1
\begin{frame}
  \titlepage
\end{frame}


% 3
\begin{frame}{Lecture Roadmap}
\begin{columns}[T]
\begin{column}{0.48\linewidth}
\textbf{Part I. What Is Data Analysis?}
\begin{itemize}
  \item Concept and objectives.
  \item Importance and applications.
  \item Main process overview.
  \item Four types of analysis.
  \item Tools and limitations.
\end{itemize}
\end{column}
\begin{column}{0.48\linewidth}
\textbf{Part II. Six-Step Process}
\begin{itemize}
  \item Define the problem.
  \item Collect data.
  \item Clean data.
  \item Analyze data.
  \item Visualize results.
  \item Interpret and decide.
\end{itemize}
\end{column}
\end{columns}
\end{frame}

\section{Part I: What Is Data Analysis?}


% 5
\begin{frame}{Learning Outcomes}
After completing this lesson, learners will be able to:
\begin{itemize}
  \item Explain the concept and objectives of data analysis.
  \item Describe the role of data analysis in decision-making.
  \item Outline the main stages of a data analysis process.
  \item Distinguish descriptive, diagnostic, predictive, and prescriptive analysis.
  \item Recognize tools such as Excel, Python, R, Tableau, Power BI, SAS, and KNIME.
  \item Describe applications in business, healthcare, finance, marketing, and scientific research.
  \item Identify limitations related to data quality, bias, resources, tools, and context.
  \item Apply the concepts to simple analytical situations.
\end{itemize}
\end{frame}


% 7
\begin{frame}{What Is Data Analysis?}
\begin{block}{Definition}
Data analysis is the process of \textbf{collecting}, \textbf{cleaning}, \textbf{transforming}, and \textbf{interpreting} data to discover useful information and support decision-making.
\end{block}
\begin{itemize}
  \item It converts raw data into meaningful information.
  \item It supports problem solving, performance evaluation, and forecasting.
  \item It is not only about software: it also requires a good question, appropriate methods, and contextual interpretation.
\end{itemize}
\end{frame}

\begin{frame}{Visual Overview: What Is Data Analysis?}
\centering
\imageplaceholder{images/image-11.png}{Data analysis overview}
\end{frame}


% 8
\begin{frame}{Data Analysis: Input, Process, Output}
\begin{center}
\begin{tabular}{p{0.22\linewidth}p{0.65\linewidth}}
\toprule
\textbf{Component} & \textbf{Meaning}\\
\midrule
Input & Raw data: possibly incomplete, inconsistent, duplicated, biased, or poorly structured.\\
Process & Cleaning, organizing, exploring, summarizing, modeling, visualizing, and interpreting.\\
Output & Insights, reports, dashboards, forecasts, recommendations, or decisions.\\
\bottomrule
\end{tabular}
\end{center}
\smallnote{Key message: data create value only when they are linked to a meaningful analytical question.}
\end{frame}

% 9
\begin{frame}{Quick Check: What Is Data Analysis?}
\begin{alertblock}{Question 1}
What is the main objective of data analysis?
\end{alertblock}
\begin{enumerate}[A.]
  \item Only to store data
  \item To transform raw data into useful information
  \item To delete all unnecessary data
  \item To completely replace humans in decision-making
\end{enumerate}
\vspace{0.5em}
\begin{alertblock}{Question 2}
Which activity is not mentioned as part of data analysis?
\end{alertblock}
\begin{enumerate}[A.]
  \item Data collection
  \item Data cleaning
  \item Data interpretation
  \item Hardware manufacturing
\end{enumerate}
\end{frame}

% 10
\begin{frame}{The Importance of Data Analysis}
Data analysis is important because it transforms raw information into insights that can be applied in practice.
\begin{itemize}
  \item \key{Data-driven decision-making}: decisions can be based on evidence rather than intuition alone.
  \item \key{Business intelligence support}: organizations can understand customers, markets, and internal operations.
  \item \key{Performance evaluation}: processes, products, services, and strategies can be measured and improved.
  \item \key{Risk management}: potential risks can be detected, forecast, and mitigated before they become serious problems.
\end{itemize}
\end{frame}

\begin{frame}{Illustration: Why Data Analysis Matters}
\centering
\imageplaceholder{images/image-14.png}{Why data analysis matters}
\end{frame}


% 11
\begin{frame}{Why Data-Driven Decisions Matter}
\begin{columns}[T]
\begin{column}{0.50\linewidth}
\textbf{Without analysis}
\begin{itemize}
  \item Decisions rely heavily on assumptions.
  \item Problems may be detected too late.
  \item The same strategy may be applied to every group.
  \item Performance is difficult to evaluate objectively.
\end{itemize}
\end{column}
\begin{column}{0.47\linewidth}
\textbf{With analysis}
\begin{itemize}
  \item Decisions are supported by evidence.
  \item Patterns and risks can be detected earlier.
  \item Actions can be customized by segment.
  \item Performance can be tracked and compared.
\end{itemize}
\end{column}
\end{columns}
\end{frame}

% 12
\begin{frame}{Quick Check: Importance (1/2)}
\begin{alertblock}{Question 1}
Why are data-driven decisions often more reliable than decisions based only on intuition?
\end{alertblock}
\begin{enumerate}[A.]
  \item Because data are always completely accurate
  \item Because data provide evidence to support decisions
  \item Because data eliminate all risks
  \item Because data do not require interpretation
\end{enumerate}
\begin{alertblock}{Question 2}
How does data analysis support risk management?
\end{alertblock}
\begin{enumerate}[A.]
  \item By completely eliminating all risks
  \item By detecting and forecasting potential risks
  \item By recording risks only after they occur
  \item By replacing all control measures
\end{enumerate}
\end{frame}

\begin{frame}{Quick Check: Importance (2/2)}
\begin{alertblock}{Question 3. True or false?}
Data analysis is valuable only for large companies.
\end{alertblock}
\end{frame}

% 13
\begin{frame}{The Data Analysis Process: Overview}
The data analysis process includes several major stages that transform raw data into valuable insights.
\begin{center}
\begin{tabular}{p{0.28\linewidth}p{0.55\linewidth}}
\toprule
\textbf{Stage} & \textbf{Core role}\\
\midrule
Define the objective & Clarify what the analysis must answer.\\
Collect data & Gather relevant and reliable data.\\
Clean and preprocess & Correct errors, missing values, duplicates, and outliers.\\
Exploratory analysis & Identify patterns, trends, relationships, and anomalies.\\
Statistical analysis & Test hypotheses, find relationships, and forecast.\\
Communicate results & Present insights for decision-making.\\
\bottomrule
\end{tabular}
\end{center}
\end{frame}

\begin{frame}{Illustration: Data Analysis Process}
\centering
\imageplaceholder{images/image-12.png}{Overview of the data analysis process}
\end{frame}


% 14
\begin{frame}{Main Stages of the Process}
\begin{enumerate}
  \item \textbf{Define the Objective}: set clear objectives and identify the main questions.
  \item \textbf{Collect Data}: collect relevant and reliable qualitative or quantitative data.
  \item \textbf{Clean and Preprocess Data}: correct errors, handle missing values, remove duplicate records, and address outliers.
  \item \textbf{Perform EDA}: use descriptive statistics and visualization to identify patterns, trends, relationships, and anomalies.
  \item \textbf{Perform Statistical Analysis}: apply statistical methods or analytical models to test hypotheses and generate forecasts.
  \item \textbf{Visualize and Communicate Results}: present findings through charts, dashboards, reports, or presentations.
\end{enumerate}
\end{frame}

% 15
\begin{frame}{Quick Check: Data Analysis Process}
\begin{alertblock}{Question 1}
What is the first stage of the data analysis process?
\end{alertblock}
\begin{enumerate}[A.]
  \item Data visualization
  \item Model building
  \item Defining the objective
  \item Data cleaning
\end{enumerate}
\begin{alertblock}{Question 2}
Handling missing values and duplicate records belongs to which stage?
\end{alertblock}
\begin{enumerate}[A.]
  \item Data collection
  \item Data cleaning and preprocessing
  \item Statistical analysis
  \item Communicating results
\end{enumerate}
\end{frame}

% 16
\begin{frame}{Quick Check: EDA and Process Order}
\begin{alertblock}{Question 3}
What is the main objective of exploratory data analysis?
\end{alertblock}
\begin{enumerate}[A.]
  \item To delete the original data
  \item To identify patterns, trends, relationships, and unusual observations
  \item To create only a written report
  \item To completely replace statistical analysis
\end{enumerate}
\vspace{0.3em}
\quiztitle{Question 4.} Arrange the following activities in a logical order:
\begin{enumerate}
  \item Clean the data
  \item Define the objective
  \item Visualize and communicate the results
  \item Collect the data
\end{enumerate}
\quiztitle{Question 5. Case.} A company wants to determine why its sales have declined. What should it do before collecting data?
\end{frame}

% 17
\begin{frame}{Types of Data Analysis}
Data analysis is commonly divided into four major types, depending on the question being addressed.
\begin{center}
\begin{tabular}{p{0.22\linewidth}p{0.25\linewidth}p{0.40\linewidth}}
\toprule
\textbf{Type} & \textbf{Question} & \textbf{Purpose}\\
\midrule
Descriptive & What happened? & Summarize historical data.\\
Diagnostic & Why did it happen? & Identify causes and contributing factors.\\
Predictive & What may happen? & Estimate future outcomes or trends.\\
Prescriptive & What should we do? & Recommend actions or strategies.\\
\bottomrule
\end{tabular}
\end{center}
\end{frame}

\begin{frame}{Illustration: Four Types of Data Analysis}
\centering
\imageplaceholder{images/image-13.png}{Four types of data analysis}
\end{frame}


% 18
\begin{frame}{Descriptive and Diagnostic Analysis}
\begin{block}{Descriptive Analysis}
Descriptive analysis summarizes historical data to explain what happened. It commonly uses reports, charts, dashboards, descriptive statistics, and performance indicators.
\end{block}
\begin{block}{Diagnostic Analysis}
Diagnostic analysis examines data in greater depth to determine why an event or outcome occurred. It helps identify possible causes, correlations, and contributing factors behind observed patterns.
\end{block}
\end{frame}

% 19
\begin{frame}{Predictive and Prescriptive Analysis}
\begin{block}{Predictive Analysis}
Predictive analysis uses historical data, statistical models, and machine-learning methods to estimate what is likely to happen in the future. It supports planning by forecasting trends, risks, and opportunities.
\end{block}
\begin{block}{Prescriptive Analysis}
Prescriptive analysis builds on predictions to recommend actions. It supports decision-making by identifying strategies or solutions likely to produce the best outcomes.
\end{block}
\end{frame}

% 20
\begin{frame}{Quick Check: Types of Data Analysis (1/2)}
\begin{enumerate}
  \item Which type of analysis answers the question ``What happened?''
  \begin{enumerate}[A.]
    \item Descriptive analysis \item Diagnostic analysis \item Predictive analysis \item Prescriptive analysis
  \end{enumerate}
  \item Which type investigates causes of an observed outcome?
  \begin{enumerate}[A.]
    \item Descriptive \item Diagnostic \item Predictive \item Prescriptive
  \end{enumerate}
\end{enumerate}
\end{frame}

\begin{frame}{Quick Check: Types of Data Analysis (2/2)}
\begin{enumerate}\setcounter{enumi}{2}
  \item Forecasting sales for the next quarter belongs to which type?
  \begin{enumerate}[A.]
    \item Descriptive \item Diagnostic \item Predictive \item Prescriptive
  \end{enumerate}
  \item A system recommends changing prices to maximize profit. Which type is this?
  \begin{enumerate}[A.]
    \item Descriptive \item Diagnostic \item Predictive \item Prescriptive
  \end{enumerate}
\end{enumerate}
\end{frame}

% 21
\begin{frame}{Matching: Core Analytical Questions}
\begin{center}
\begin{tabular}{p{0.42\linewidth}p{0.42\linewidth}}
\toprule
\textbf{Question} & \textbf{Type of analysis}\\
\midrule
What happened? & Descriptive analysis\\
Why did it happen? & Diagnostic analysis\\
What may happen? & Predictive analysis\\
What should we do? & Prescriptive analysis\\
\bottomrule
\end{tabular}
\end{center}
\end{frame}

% 22
\begin{frame}{Data Analysis Tools}
\begin{center}\scriptsize
\begin{tabular}{p{0.19\linewidth}p{0.70\linewidth}}
\toprule
\textbf{Tool} & \textbf{Main use}\\
\midrule
Microsoft Excel & Basic calculations, pivot tables, data summaries, and simple charts.\\
SAS & Advanced analysis, statistical modeling, and predictive analytics.\\
R & Statistical analysis, data modeling, and visualization.\\
Python & Data processing, analysis, visualization, and machine learning with Pandas, NumPy, and Matplotlib.\\
Tableau & Interactive dashboards and advanced data visualization.\\
Power BI & Business-intelligence reports and dashboards in the Microsoft ecosystem.\\
KNIME & Open-source workflows for data mining, processing, and machine learning.\\
\bottomrule
\end{tabular}
\end{center}
\end{frame}

% 23
\begin{frame}{Quick Check: Tools}
\begin{enumerate}
  \item Which tool is suitable for basic calculations, pivot tables, and simple charts?
  \begin{enumerate}[A.]
    \item Excel \item KNIME \item SAS \item Tableau
  \end{enumerate}
  \item Pandas, NumPy, and Matplotlib are commonly used with which language?
  \begin{enumerate}[A.]
    \item R \item Python \item SQL \item Java
  \end{enumerate}
  \item Which tools are commonly used to build interactive dashboards?
  \begin{enumerate}[A.]
    \item Tableau or Power BI \item Only a word processor \item A file browser \item A pocket calculator
  \end{enumerate}
\end{enumerate}
\quiztitle{Question 4. True or false?} Every data-analysis problem must use the same tool.
\end{frame}

% 24
\begin{frame}{Applications of Data Analysis}
\begin{itemize}
  \item \key{Business intelligence}: monitor sales, customer behavior, operational performance, and strategy.
  \item \key{Healthcare}: monitor treatment outcomes, evaluate procedures, improve operations, and support medical research.
  \item \key{Finance}: detect fraud, evaluate risk, analyze investments, forecast markets, and support budgeting.
  \item \key{Marketing}: understand customer preferences, evaluate campaigns, segment markets, and design targeted strategies.
  \item \key{Scientific research}: interpret experimental results, test hypotheses, identify patterns, and generate knowledge.
\end{itemize}
\end{frame}

% 25
\begin{frame}{Quick Check: Applications}
\begin{enumerate}
  \item Fraud detection is a common application of data analysis in which field?
  \begin{enumerate}[A.]
    \item Finance \item Linguistics \item Graphic design \item Architecture
  \end{enumerate}
  \item Customer segmentation and campaign-performance evaluation belong to which field?
  \begin{enumerate}[A.]
    \item Healthcare \item Marketing \item Mechanical engineering \item Physics
  \end{enumerate}
  \item In healthcare, data analysis can be used to:
  \begin{enumerate}[A.]
    \item Monitor treatment outcomes
    \item Evaluate treatment effectiveness
    \item Support medical research
    \item All of the above
  \end{enumerate}
\end{enumerate}
\quiztitle{Question 4. Case.} A university wants to identify factors associated with student dropout risk. In which application area does this problem belong?
\end{frame}

% 26
\begin{frame}{Limitations of Data Analysis}
\begin{itemize}
  \item \key{Data-quality problems}: inaccurate, incomplete, outdated, or inconsistent data can mislead.
  \item \key{Time and resource requirements}: collection, cleaning, processing, and interpretation require effort.
  \item \key{Risk of bias}: biased datasets, methods, or assumptions may produce unreliable results.
  \item \key{Overreliance on tools}: software output can be wrong or misleading without methodological knowledge.
  \item \key{Lack of context}: quantitative data may not capture qualitative, social, behavioral, or human factors.
\end{itemize}
\end{frame}

% 27
\begin{frame}{Quick Check: Limitations}
\begin{enumerate}
  \item What may happen when data are incomplete or inconsistent?
  \begin{enumerate}[A.]
    \item Results always become more accurate
    \item Conclusions may be biased or misleading
    \item Data cleaning becomes unnecessary
    \item The model automatically corrects every error
  \end{enumerate}
  \item Why should analysts not rely entirely on analytical tools?
  \begin{enumerate}[A.]
    \item Because every tool is inaccurate
    \item Because results must be evaluated using methodological knowledge and context
    \item Because tools can only be used with small datasets
    \item Because tools cannot create charts
  \end{enumerate}
\end{enumerate}
\quiztitle{Question 3. True or false?} A highly accurate result obtained from biased data may still lead to an inappropriate decision.\\[0.3em]
\quiztitle{Question 4.} What does the limitation ``lack of context'' mean?
\end{frame}

% 28
\begin{frame}{Conclusion: What Data Analysis Provides}
\begin{block}{Summary}
Data analysis provides a systematic method for transforming raw data into useful knowledge.
\end{block}
\begin{itemize}
  \item Define objectives.
  \item Collect and prepare data.
  \item Explore patterns and apply analytical methods.
  \item Communicate results for decision-making.
  \item Evaluate reliability through data quality, method suitability, and correct interpretation.
\end{itemize}
\end{frame}

% 29
\begin{frame}{End-of-Lesson Review: Part A}
\begin{enumerate}\scriptsize
  \item Data analysis is the process of:
  \begin{enumerate}[A.]\scriptsize
    \item Only collecting and storing data
    \item Collecting, cleaning, transforming, and interpreting data
    \item Only visualizing data
    \item Only building machine-learning models
  \end{enumerate}
  \item What should be done before collecting data?
  \begin{enumerate}[A.]\scriptsize
    \item Choose chart colors
    \item Define the objective and analytical questions
    \item Build the final report
    \item Remove all outliers
  \end{enumerate}
  \item Which type of analysis answers ``Why did sales decline?''
  \begin{enumerate}[A.]\scriptsize
    \item Descriptive \item Diagnostic \item Predictive \item Prescriptive
  \end{enumerate}
  \item Which type answers ``What might next month's sales be?''
  \begin{enumerate}[A.]\scriptsize
    \item Descriptive \item Diagnostic \item Predictive \item Prescriptive
  \end{enumerate}
\end{enumerate}
\end{frame}

% 30
\begin{frame}{End-of-Lesson Review: Part A Continued}
\begin{enumerate}\scriptsize\setcounter{enumi}{4}
  \item Which type of analysis answers ``Which strategy should the company choose?''
  \begin{enumerate}[A.]\scriptsize
    \item Descriptive \item Diagnostic \item Predictive \item Prescriptive
  \end{enumerate}
  \item Which of the following is a programming language widely used in data science?
  \begin{enumerate}[A.]\scriptsize
    \item Python \item Tableau \item Power BI \item Excel
  \end{enumerate}
  \item Missing values are usually handled during which stage?
  \begin{enumerate}[A.]\scriptsize
    \item Defining the objective \item Data cleaning and preprocessing \item Communicating results \item Prescriptive analysis
  \end{enumerate}
  \item Which is a limitation of data analysis?
  \begin{enumerate}[A.]\scriptsize
    \item Data always capture the full context
    \item No expertise is needed to interpret results
    \item Biased data can produce unreliable conclusions
    \item Data analysis eliminates all risk
  \end{enumerate}
\end{enumerate}
\end{frame}

% 31
\begin{frame}{End-of-Lesson Review: True/False and Short Answers}
\begin{block}{Part B. True/False}
\begin{enumerate}
  \item Descriptive analysis focuses on explaining what happened.
  \item Predictive analysis guarantees that the future will occur exactly as forecast.
  \item Visualization helps communicate analytical findings to stakeholders.
  \item Poor-quality data do not affect the results if powerful software is used.
  \item Contextual knowledge is important when interpreting results.
\end{enumerate}
\end{block}
\begin{block}{Part C. Short-answer prompts}
\begin{enumerate}
  \item Present the six major stages of the data analysis process.
  \item Distinguish between descriptive and diagnostic analysis.
  \item Why does data quality strongly affect analytical results?
  \item State two benefits and two limitations of data analysis.
\end{enumerate}
\end{block}
\end{frame}

% 32
\begin{frame}{Case-Based Exercises}
\begin{block}{Case 1: Sales Analysis}
A store observes that June sales are 15\% lower than May sales.
\begin{enumerate}
  \item What information should descriptive analysis provide?
  \item Which factors should diagnostic analysis investigate?
  \item How can predictive analysis be used?
  \item What recommendations might prescriptive analysis provide?
\end{enumerate}
\end{block}
\begin{block}{Case 2: Data Quality}
A customer dataset contains many duplicate rows, missing age values, and several negative income values.
\begin{enumerate}
  \item Which data-quality problems are present?
  \item During which stage should these problems be handled?
  \item What may happen if the dataset is used without cleaning?
\end{enumerate}
\end{block}
\end{frame}

\section{Part II: Six Steps in the Data Analysis Process}


% 34
\begin{frame}{Learning Outcomes}
After completing this lesson, learners will be able to:
\begin{itemize}\small
  \item Explain the role of a structured data analysis process.
  \item Define the problem, objectives, and success criteria of an analytical task.
  \item Identify and select appropriate data sources.
  \item Examine the structure, origin, and initial quality of a dataset.
  \item Handle missing values, unnecessary columns, and categorical variables.
  \item Use Python, Pandas, Seaborn, and Matplotlib to analyze and visualize data.
  \item Read and interpret correlation matrices, bar charts, histograms, and scatter plots.
  \item Understand basic data splitting, model training, and accuracy evaluation.
  \item Convert analytical results into recommendations and decisions.
  \item Recognize that correlation does not imply causation and that a single metric is not sufficient to judge a model.
\end{itemize}
\end{frame}


% 36
\begin{frame}[fragile]{Installing Required Libraries}
\begin{lstlisting}[style=py]
pip install pandas seaborn matplotlib scikit-learn
\end{lstlisting}
\begin{itemize}
  \item \texttt{pandas}: tabular data manipulation.
  \item \texttt{seaborn}: statistical visualization and built-in datasets.
  \item \texttt{matplotlib}: plotting backend and chart customization.
  \item \texttt{scikit-learn}: train-validation split, models, and evaluation metrics.
\end{itemize}
\end{frame}

% 37
\begin{frame}{Why Use a Structured Process?}
Data analysis follows a structured approach in which:
\begin{itemize}
  \item \key{Step-by-step process}: raw data are transformed into meaningful insights.
  \item \key{Systematic approach}: the process helps improve accuracy and reliability.
  \item \key{Better decision-making}: decisions are based on evidence rather than intuition alone.
\end{itemize}
\begin{alertblock}{Quick Check}
\begin{enumerate}
  \item What is the main objective of a structured data analysis process?
  \item True or false? A systematic data analysis process helps improve the accuracy and reliability of results.
\end{enumerate}
\end{alertblock}
\end{frame}

% 38
\begin{frame}{The Six-Step Workflow}
\begin{center}
\begin{tabular}{p{0.28\linewidth}p{0.56\linewidth}}
\toprule
\textbf{Step} & \textbf{Central question}\\
\midrule
Define the problem & What problem must be solved?\\
Collect data & What data are needed and where can they be obtained?\\
Clean the data & Are the data ready for analysis?\\
Analyze the data & What do the data show?\\
Visualize the results & How can the results be made easier to understand?\\
Interpret and decide & What do the results mean and what should be done next?\\
\bottomrule
\end{tabular}
\end{center}
\end{frame}

\begin{frame}{Illustration: Six-Step Workflow}
\centering
\imageplaceholder{images/image-4-v1.png}{Structured data analysis workflow}
\end{frame}


% 39
\begin{frame}{Step 1: Define the Problem}
Before beginning any analytical activity, the problem must be clearly defined.
\begin{itemize}
  \item Clarify the question, objective, or analytical opportunity.
  \item Ensure the analytical goal aligns with stakeholder expectations.
  \item Keep the analysis focused and relevant.
  \item Avoid unnecessary data collection.
\end{itemize}
\begin{block}{Main tasks}
Identify the core problem, define expected outcomes, understand context and constraints, and define success criteria.
\end{block}
\end{frame}

% 40
\begin{frame}{Step 1 Example: Sales Decline}
A company notices that sales have declined significantly over the last three months.
\begin{block}{Instead of asking only ``Why did sales decline?'', ask:}
\begin{itemize}
  \item Which products experienced the largest decline?
  \item Which region had the greatest decrease?
  \item Which customer groups were most affected?
  \item Is the decline related to pricing, inventory, or marketing activities?
\end{itemize}
\end{block}
\end{frame}

% 41
\begin{frame}{Quick Check: Step 1}
\begin{enumerate}
  \item Why should the problem be defined before data collection?
  \begin{enumerate}[A.]
    \item To reduce file size
    \item To keep the analysis focused and relevant
    \item To avoid using charts
    \item To remove all qualitative data
  \end{enumerate}
  \item Which activity belongs to the problem-definition step?
  \begin{enumerate}[A.]
    \item Handling missing values
    \item Defining success criteria
    \item Training a model
    \item Drawing a heatmap
  \end{enumerate}
\end{enumerate}
\quiztitle{Question 3. Case.} A university wants to investigate why students drop out. Suggest two more specific analytical questions.
\end{frame}

% 42
\begin{frame}{Step 2: Collect Data}
After defining the problem, collect data from appropriate sources.
\begin{itemize}
  \item Data may come from internal databases, APIs, surveys, web scraping, or public datasets.
  \item Data must be relevant, accurate, and sufficiently complete.
  \item Combining multiple sources may enrich the analysis.
  \item The origin and structure of each dataset should be documented for transparency.
  \item Update frequency, file format, and refresh requirements should be considered.
\end{itemize}
\end{frame}

% 43
\begin{frame}[fragile]{Step 2 Example: Titanic Dataset}
\begin{lstlisting}[style=py]
import seaborn as sns
import pandas as pd

titanic = sns.load_dataset("titanic")
titanic.head()
\end{lstlisting}
\begin{itemize}
  \item The Titanic dataset is built into Seaborn.
  \item \texttt{head()} displays the first few rows.
  \item The dataset is used to demonstrate cleaning, analysis, visualization, and model building.
\end{itemize}
\end{frame}

\begin{frame}{Illustration: Titanic Dataset Preview}
\centering
\imageplaceholder{images/image-5-v1.png}{Titanic dataset preview}
\end{frame}


% 44
\begin{frame}{Quick Check: Step 2}
\begin{enumerate}
  \item Which of the following can be used as a data source?
  \begin{enumerate}[A.]
    \item Internal databases \item APIs \item Surveys \item All of the above
  \end{enumerate}
  \item Why should the origin of a dataset be documented?
  \begin{enumerate}[A.]
    \item To improve transparency and auditability
    \item To make the dataset larger
    \item To avoid data cleaning
    \item To replace data analysis
  \end{enumerate}
  \item Which function is used to load the Titanic dataset in the example?
  \begin{enumerate}[A.]
    \item \texttt{pd.read\_csv()} \item \texttt{sns.load\_dataset()} \item \texttt{plt.load()} \item \texttt{np.loadtxt()}
  \end{enumerate}
\end{enumerate}
\end{frame}

% 45
\begin{frame}{Step 3: Clean the Data}
Raw data are rarely ready for direct analysis.
\begin{block}{Data cleaning includes}
\begin{itemize}
  \item Handling missing values.
  \item Removing duplicate records.
  \item Standardizing formats.
  \item Converting categorical variables into numerical form.
  \item Removing irrelevant, redundant, or inconsistent columns.
\end{itemize}
\end{block}
Well-prepared data improve the reliability and accuracy of analytical results.
\end{frame}

% 46
\begin{frame}[fragile]{Step 3: Check and Handle Missing Data}
\begin{lstlisting}[style=py]
print(titanic.isnull().sum())

titanic["age"].fillna(
    titanic["age"].median(),
    inplace=True
)

titanic["embarked"].fillna(
    titanic["embarked"].mode()[0],
    inplace=True
)
\end{lstlisting}
\begin{itemize}
  \item \texttt{isnull().sum()} counts missing values in each column.
  \item Missing values in \texttt{age} are replaced by the median.
  \item Missing values in \texttt{embarked} are replaced by the most frequent value.
\end{itemize}
\end{frame}

% 47
\begin{frame}[fragile]{Step 3: Remove Columns and Encode Categories}
\begin{lstlisting}[style=py]
titanic.drop(
    ["deck", "embark_town", "alive", "class", "who", "adult_male"],
    axis=1,
    inplace=True
)

titanic["sex"] = titanic["sex"].map({"male": 0, "female": 1})
titanic["embarked"] = titanic["embarked"].map({"C": 0, "Q": 1, "S": 2})

titanic.head()
\end{lstlisting}
\smallnote{Categorical variables often need to be converted before analysis or modeling.}
\end{frame}

\begin{frame}{Illustration: Cleaned Titanic Data}
\centering
\imageplaceholder{images/image-6.png}{Titanic data after cleaning and encoding}
\end{frame}


% 48
\begin{frame}{Quick Check: Step 3}
\begin{enumerate}\small
  \item Which command counts missing values in each column?
  \begin{enumerate}[A.]\small
    \item \texttt{titanic.head()} \item \texttt{titanic.isnull().sum()} \item \texttt{titanic.describe()} \item \texttt{titanic.drop()}
  \end{enumerate}
  \item In the example, missing values in \texttt{age} are replaced by:
  \begin{enumerate}[A.]\small
    \item The minimum value \item The maximum value \item The median \item Zero
  \end{enumerate}
  \item Why is the \texttt{sex} column converted into numerical values?
  \begin{enumerate}[A.]\small
    \item To reduce rows \item To make the data more suitable for analytical algorithms or models \item To change gender \item To create a pie chart
  \end{enumerate}
\end{enumerate}
\quiztitle{Question 4. True or false?} Every column containing missing values can be removed without considering context.
\end{frame}

% 49
\begin{frame}{Step 4: Analyze the Data}
Data analysis is the core step in which the analyst searches for patterns, trends, and relationships.
\begin{itemize}
  \item Calculate measures such as mean, median, mode, and variance.
  \item Identify correlations, trends, and unusual observations.
  \item Apply models such as regression, clustering, or classification.
  \item Compare results with expectations or hypotheses.
\end{itemize}
\end{frame}

% 50
\begin{frame}[fragile]{Step 4: Correlation Matrix}
\begin{lstlisting}[style=py]
import matplotlib.pyplot as plt

plt.figure(figsize=(8, 6))
sns.heatmap(
    titanic.corr(),
    annot=True,
    cmap="coolwarm"
)

plt.title("Correlation Matrix")
plt.show()
\end{lstlisting}
A heatmap helps observe the degree of correlation among numerical variables.
\end{frame}

\begin{frame}{Illustration: Correlation Matrix}
\centering
\imageplaceholder{images/image-7.png}{Correlation matrix for the Titanic dataset}
\end{frame}


% 51
\begin{frame}[fragile]{Step 4: Survival Rate by Passenger Class}
\begin{lstlisting}[style=py]
sns.barplot(
    x="pclass",
    y="survived",
    data=titanic
)

plt.title("Survival Rate by Passenger Class")
plt.show()
\end{lstlisting}
The bar chart compares the average survival rate across passenger classes.
\end{frame}

\begin{frame}{Illustration: Survival Rate by Passenger Class}
\centering
\imageplaceholder{images/image-8.png}{Survival rate by passenger class}
\end{frame}


% 52
\begin{frame}{Quick Check: Step 4}
\begin{enumerate}
  \item What is a correlation matrix used for?
  \begin{enumerate}[A.]
    \item Checking file types \item Examining relationships among variables \item Removing duplicates \item Loading data from an API
  \end{enumerate}
  \item In the example, the barplot compares:
  \begin{enumerate}[A.]
    \item Age and fare \item Passenger class and survival rate \item Gender and port \item Number of columns and rows
  \end{enumerate}
  \item Which models may be used during data analysis?
  \begin{enumerate}[A.]
    \item Regression \item Clustering \item Classification \item All of the above
  \end{enumerate}
\end{enumerate}
\quiztitle{Question 4. Case.} If two variables are highly correlated, can we conclude with certainty that one causes the other?
\end{frame}

% 53
\begin{frame}{Step 5: Visualize the Results}
Visualization makes complex data easier to understand.
\begin{itemize}
  \item Charts, graphs, and dashboards highlight important trends, patterns, and unusual observations.
  \item A good visualization should be clear, intuitive, and useful for decision-making.
  \item Select chart types according to analytical questions: histogram, scatter plot, bar chart, heatmap, etc.
  \item Avoid excessive detail that distracts from the main message.
\end{itemize}
\end{frame}

% 54
\begin{frame}[fragile]{Step 5: Count Plot and Age Distribution}
\begin{lstlisting}[style=py]
sns.countplot(
    x="survived",
    data=titanic
)
plt.title("Number of Passengers by Survival Status")
plt.show()

sns.histplot(
    titanic["age"],
    kde=True
)
plt.title("Age Distribution")
plt.show()
\end{lstlisting}
\end{frame}


\begin{frame}{Illustration: Survival Count and Age Distribution}
\begin{columns}[c]
\begin{column}{0.50\linewidth}
\centering
\imageplaceholder{images/image-9.png}{Passenger counts by survival status}
\end{column}
\begin{column}{0.50\linewidth}
\centering
\imageplaceholder{images/image-1-v1.png}{Age distribution of Titanic passengers}
\end{column}
\end{columns}
\end{frame}


% 55
\begin{frame}[fragile]{Step 5: Scatter Plot of Age and Fare}
\begin{lstlisting}[style=py]
sns.scatterplot(
    x="age",
    y="fare",
    hue="survived",
    data=titanic
)

plt.title("Relationship Between Fare and Age by Survival Status")
plt.show()
\end{lstlisting}
\smallnote{The \texttt{hue} argument colors points by survival status.}
\end{frame}

\begin{frame}{Illustration: Fare vs Age by Survival Status}
\centering
\imageplaceholder{images/image-10.png}{Age and fare by survival status}
\end{frame}


% 56
\begin{frame}{Quick Check: Step 5}
\begin{enumerate}
  \item Which chart is suitable for displaying the distribution of \texttt{age}?
  \begin{enumerate}[A.]
    \item Histogram \item Network chart \item Gantt chart \item Geographic map
  \end{enumerate}
  \item What is the purpose of \texttt{hue="survived"} in the scatter plot?
  \begin{enumerate}[A.]
    \item To remove the \texttt{survived} column
    \item To distinguish points by survival status
    \item To change the figure size
    \item To calculate the mean
  \end{enumerate}
\end{enumerate}
\quiztitle{Question 3. True or false?} A chart is always more effective when it contains more details.\\[0.4em]
\quiztitle{Question 4.} Why should the chart type be selected according to the analytical question?
\end{frame}

% 57
\begin{frame}{Step 6: Interpret Results and Make Decisions}
The final step is to transform analytical results into actionable insights.
\begin{itemize}
  \item Place results in the context of the original problem.
  \item Provide actionable recommendations based on findings.
  \item Communicate results clearly to stakeholders.
  \item Monitor outcomes and repeat the process for continuous improvement.
\end{itemize}
\smallnote{Interpretation connects the numerical result with practical decision-making.}
\end{frame}

% 58
\begin{frame}[fragile]{Step 6: Building a Survival Prediction Model}
\begin{lstlisting}[style=py]
X = titanic.drop("survived", axis=1)
y = titanic["survived"]

X_train, X_val, y_train, y_val = train_test_split(
    X, y, test_size=0.2, random_state=42
)

model = RandomForestClassifier(n_estimators=100, random_state=42)
model.fit(X_train, y_train)

y_pred = model.predict(X_val)
accuracy = accuracy_score(y_val, y_pred)
print(f"Model accuracy: {accuracy:.4f}")
\end{lstlisting}
\end{frame}

% 59
\begin{frame}{Step 6: Interpreting Accuracy}
\begin{block}{Result}
\texttt{Model accuracy: 0.8101}
\end{block}
\begin{itemize}
  \item The model correctly predicts approximately 81.01\% of observations in the validation set.
  \item Accuracy alone should not be used to evaluate a model.
  \item Practical applications may require additional metrics and careful context-specific interpretation.
  \item The code uses \texttt{train\_test\_split}, \texttt{RandomForestClassifier}, and \texttt{accuracy\_score}, which must be imported from Scikit-learn.
\end{itemize}
\end{frame}

% 60
\begin{frame}{Quick Check: Step 6}
\begin{enumerate}\small
  \item What is the objective of the interpretation step?
  \begin{enumerate}[A.]\small
    \item Only to present numbers
    \item To transform analytical results into insights and actions
    \item To delete all data
    \item To change the original problem
  \end{enumerate}
  \item An accuracy value of \texttt{0.8101} is approximately:
  \begin{enumerate}[A.]\small
    \item 8.101\% \item 18.01\% \item 81.01\% \item 810.1\%
  \end{enumerate}
  \item Why should results be monitored after a decision is implemented?
  \begin{enumerate}[A.]\small
    \item To determine whether the decision is effective and make adjustments when needed
    \item To increase the number of columns
    \item To avoid communicating with stakeholders
    \item To ensure no further analysis is needed
  \end{enumerate}
\end{enumerate}
\quiztitle{Question 4. True or false?} Accuracy alone is always sufficient to fully evaluate a classification model.
\end{frame}

% 61
\begin{frame}{Process Summary}
\begin{center}\small
\begin{tabular}{p{0.27\linewidth}p{0.34\linewidth}p{0.27\linewidth}}
\toprule
\textbf{Step} & \textbf{Main content} & \textbf{Central question}\\
\midrule
Define the problem & Objective, context, success criteria & What problem must be solved?\\
Collect data & Relevant sources & What data are needed?\\
Clean data & Missing values, formatting & Are data ready?\\
Analyze data & Patterns and relationships & What do data show?\\
Visualize results & Charts and dashboards & How to make results clear?\\
Interpret and decide & Recommendations and actions & What should be done?\\
\bottomrule
\end{tabular}
\end{center}
\end{frame}

% 62
\begin{frame}{Review: Multiple-Choice Questions 1-4}
\begin{enumerate}\scriptsize
  \item What is the first step in the data analysis process?
  \begin{enumerate}[A.]\scriptsize
    \item Clean the data \item Define the problem \item Visualize the data \item Train a model
  \end{enumerate}
  \item Which activity belongs to the data-collection step?
  \begin{enumerate}[A.]\scriptsize
    \item Identify data sources \item Handle missing values \item Calculate a correlation matrix \item Evaluate model accuracy
  \end{enumerate}
  \item The median is commonly used to:
  \begin{enumerate}[A.]\scriptsize
    \item Rename a column \item Fill missing values in a numerical variable \item Create an API \item Draw a scatter plot
  \end{enumerate}
  \item Which function is used to visualize a correlation matrix in the example?
  \begin{enumerate}[A.]\scriptsize
    \item \texttt{sns.heatmap()} \item \texttt{sns.load\_dataset()} \item \texttt{pd.DataFrame()} \item \texttt{model.fit()}
  \end{enumerate}
\end{enumerate}
\end{frame}

% 63
\begin{frame}{Review: Multiple-Choice Questions 5-8}
\begin{enumerate}\scriptsize\setcounter{enumi}{4}
  \item Which chart is used to show the age distribution?
  \begin{enumerate}[A.]\scriptsize
    \item Histogram \item Pie chart \item Gantt chart \item Geographic heatmap
  \end{enumerate}
  \item What does \texttt{test\_size=0.2} mean?
  \begin{enumerate}[A.]\scriptsize
    \item 20\% of data are used as validation \item 20\% are deleted \item There are 20 input variables \item Model has 20\% accuracy
  \end{enumerate}
  \item An accuracy of \texttt{0.8101} is equal to:
  \begin{enumerate}[A.]\scriptsize
    \item 0.8101\% \item 8.101\% \item 81.01\% \item 810.1\%
  \end{enumerate}
  \item Which step converts analytical results into recommendations?
  \begin{enumerate}[A.]\scriptsize
    \item Data collection \item Data cleaning \item Interpretation and decision-making \item Visualization only
  \end{enumerate}
\end{enumerate}
\end{frame}

% 64
\begin{frame}{Review: True/False}
\begin{block}{Part B. True/False}
\begin{enumerate}
  \item The analytical objective should be defined before data collection.
  \item Raw data can always be used directly to build a model.
  \item A correlation matrix can help identify relationships among variables.
  \item A high correlation always proves causation.
  \item Data visualization can support communication with stakeholders.
  \item After a decision has been made, there is no need to monitor the outcome.
\end{enumerate}
\end{block}
\end{frame}

\begin{frame}{Review: Short Answers}
\begin{block}{Part C. Short answers}
\begin{enumerate}
  \item Present the six main steps.
  \item Why is defining the problem important?
  \item State three common problems in raw data.
  \item Distinguish data analysis and visualization.
  \item Why must model results be interpreted in context?
\end{enumerate}
\end{block}
\end{frame}

% 65
\begin{frame}{Practical Exercises}
\begin{block}{Exercise 1. Explore the Titanic Dataset}
Display the first five rows, determine rows and columns, check data types, count missing values, and comment on quality.
\end{block}
\begin{block}{Exercise 2. Clean the Data}
Fill missing values, remove unnecessary columns, convert \texttt{sex} into numerical values, and recheck the data.
\end{block}
\begin{block}{Exercise 3. Analyze and Visualize}
Calculate descriptive statistics, plot age distribution, compare survival by class, draw age/fare scatter plot, and write observations.
\end{block}
\end{frame}

% 66
\begin{frame}{Practical Exercise 4: Interpret the Result}
Assume that a survival-prediction model achieves an accuracy of 81.01\%.
\begin{enumerate}
  \item Explain what the result means.
  \item State two reasons why accuracy alone should not be used.
  \item Suggest two additional evaluation metrics.
  \item State one ethical or practical limitation of using the model.
\end{enumerate}
\end{frame}

\appendix
\section{Answer Key}

% 67
\begin{frame}{Answer Key: What Is Data Analysis?}
\begin{itemize}\small
  \item Quick Check - What Is Data Analysis: Q1 B; Q2 D.
  \item Importance: Q1 B; Q2 B; Q3 False.
  \item Data Analysis Process: Q1 C; Q2 B; Q3 B; Q4 order is 2 $\rightarrow$ 4 $\rightarrow$ 1 $\rightarrow$ 3.
  \item Types: Q1 A; Q2 B; Q3 C; Q4 D.
  \item Tools: Q1 A; Q2 B; Q3 A; Q4 False.
  \item Applications: Q1 A; Q2 B; Q3 D; Q4 Education.
  \item Limitations: Q1 B; Q2 B; Q3 True; Q4 quantitative data may miss qualitative, social, behavioral, or human context.
\end{itemize}
\end{frame}

% 68
\begin{frame}{Answer Key: Review for Part I}
\begin{itemize}\small
  \item Multiple choice: 1 B; 2 B; 3 B; 4 C; 5 D; 6 A; 7 B; 8 C.
  \item True/False: 1 True; 2 False; 3 True; 4 False; 5 True.
  \item Short answer Q1: define objective, collect data, clean/preprocess data, perform EDA, perform statistical analysis, visualize/communicate results.
  \item Short answer Q2: descriptive explains what happened; diagnostic investigates why it happened.
  \item Short answer Q3: poor-quality data can produce misleading patterns, biased conclusions, and unreliable models.
  \item Short answer Q4: benefits include better decisions and risk management; limitations include bias, poor quality, resource needs, and lack of context.
\end{itemize}
\end{frame}

% 69
\begin{frame}{Answer Key: Six-Step Process Quick Checks}
\begin{itemize}\small
  \item Introduction: Q1 B; Q2 True.
  \item Step 1: Q1 B; Q2 B; case examples: student groups with highest dropout rates; relationship with grades, finances, or engagement.
  \item Step 2: Q1 D; Q2 A; Q3 B.
  \item Step 3: Q1 B; Q2 C; Q3 B; Q4 False.
  \item Step 4: Q1 B; Q2 B; Q3 D; Q4 No - correlation does not imply causation.
  \item Step 5: Q1 A; Q2 B; Q3 False; Q4 different chart types fit different data and communication objectives.
  \item Step 6: Q1 B; Q2 C; Q3 A; Q4 False.
\end{itemize}
\end{frame}

% 70
\begin{frame}{Answer Key: Review for Part II}
\begin{itemize}\small
  \item Multiple choice: 1 B; 2 A; 3 B; 4 A; 5 A; 6 A; 7 C; 8 C.
  \item True/False: 1 True; 2 False; 3 True; 4 False; 5 True; 6 False.
  \item Short answer Q1: define problem, collect data, clean data, analyze data, visualize results, interpret and decide.
  \item Q2: defining the problem determines the objective, scope, required data, and success criteria.
  \item Q3: common raw-data problems include missing values, duplicates, inconsistent formats, outliers, and unsuitable data types.
  \item Q4: analysis discovers patterns and draws conclusions; visualization presents data or findings visually.
  \item Q5: context determines what a statistical or model result actually means for decisions.
\end{itemize}
\end{frame}

% 71
\begin{frame}{Closing Message}
\begin{block}{Core idea}
Good data analysis is not a single tool or chart. It is a disciplined process that starts with a clear question, prepares data carefully, explores patterns, communicates findings, and connects results to decisions.
\end{block}
\begin{itemize}
  \item Use the process to avoid jumping directly into tools.
  \item Use quizzes to check conceptual understanding.
  \item Use exercises to build practical fluency with real datasets.
\end{itemize}
\end{frame}

\end{document}
